\documentclass[12pt]{article}
\usepackage[utf8]{inputenc}
\usepackage[T1]{fontenc}
\usepackage{hyperref}
\title{Vo Wa Language Resources for K-12 Curriculum Development}
\author{Compiled by Claude Sonnet 5 (claude-sonnet-5) for Andr\'as Kornai}
\date{}
\begin{document}
\maketitle
\section{Introduction}
Vo Wa (ISO 639-3 wbm) is a variety of Wa associated with Zhenkang, China. The searches returned material on Wa generally, not Vo specifically. What follows is Wa material that may be relevant, and its applicability to Vo is not established.
\section{Primary Educational Resources}
A bilingual Chinese--Wa primary-school text, \emph{L\=ai Loux} (2003), is reported in the Wikipedia article (not seen directly; variety not specified): \url{https://en.wikipedia.org/wiki/Wa_language}
\section{Linguistic and Reference Materials}
\begin{itemize}
\item Watkins, J. (2013), \emph{Dictionary of Wa} (2 vols), Brill; about 12,000 headwords glossed in Burmese, Chinese and English, with headwords also in the China-based spelling: \url{https://eprints.soas.ac.uk/17317}
\item SOAS Wa Dictionary Project: \url{https://humancomp.org/wadict/index.html}
\item Wa resources page: \url{https://humancomp.org/wadict/wa_resources.html}
\item Wa corpus page: \url{https://humancomp.org/wadict/wa_corpus.html}
\item Omniglot on Wa writing (Latin orthographies, including a roman system developed in China in the 1950s): \url{https://www.omniglot.com/writing/wa.htm}
\end{itemize}
\section{Religious and Community Resources}
A Bible-based Latin orthography dates to the 1930s (per search summaries); no specific scripture URL was found.
\section{Digital and Audio Resources}
See the SOAS corpus page above.
\section{Curriculum Development Context}
Earlier Wa--Chinese dictionaries (1981, 1987) are described as limited or out of date. Which orthography fits Vo Wa speakers in China needs local confirmation.
\end{document}